\usepackage{booktabs}
\usepackage{amsmath}
\usepackage{amsthm}
%\usepackage[T1]{fontenc}
%\usepackage[utf8]{inputenc}
\usepackage{microtype}
\usepackage{multicol}
\usepackage{xcolor}
\usepackage{fontspec}
\usepackage{float}
%\setmainfont{Arial} % possibile opzione
\setmainfont{Palatino} % per Mac
%\setmainfont{Palatino Linotype} % per Windows

\usepackage{graphicx}
\usepackage{geometry}
\usepackage{hyperref}
\usepackage{setspace}
%\onehalfspacing
\singlespacing

\makeatletter
\def\thm@space@setup{%
  \thm@preskip=8pt plus 2pt minus 4pt
  \thm@postskip=\thm@preskip
}
\makeatother

\addto\captionsitalian{%
  \renewcommand{\chaptername}{}%
}

\usepackage{titlesec}
\titleformat{\chapter}[hang]{\huge\bfseries}{\thechapter\quad}{0pt}{}



% Colore istituzionale Sapienza
\definecolor{sapred}{rgb}{0.5098039,0.1411765,0.2}

\hypersetup{
  colorlinks=true,
  linkcolor=sapred,
  urlcolor=sapred,
  pdftitle={Titolo della tesi},
  pdfauthor={Nome Cognome}
}